\chapter[Дискретная свертка. Реализация приема и передачи BPSK сигналов]{Дискретная свертка. Реализация приема\\ и передачи BPSK сигналов}

\begin{leftbar}
    \href{https://github.com/someotb/sdr_reports/tree/main/practices/practice5}{Ссылка на GitHub}
\end{leftbar}

\section*{Цель практики:}

Програмно реализовать дискретную свертку, используя язык программировани C++,
передать полученные после свертки биты, посмотреть, что мы получим на выходе

\section*{Краткие теоретические сведения}

Дискретная свертка — это математическая операция над двумя дискретными последовательностями (сигналами), определяемая следующим образом:

Пусть даны две последовательности x[n] и h[n], где n $n\in{Z}$. Их свертка y[n] вычисляется по формуле:

\vspace{20pt}

$y[n] = (x * h)[n] = \sum_{k=-\infty}^{\infty} x[k] \cdot h[n - k]$

\vspace{20pt}

Для последовательностей конечной длины L и M, где x[n] определена при 0 $\leq n \leq L-1$, а  h[n] — при 0 $\leq n \leq M-1$, результирующая последовательность \( y[n] \) будет иметь длину \( N = L + M - 1 \) и формула принимает вид:

\vspace{20pt}

$y[n] = \sum_{k=\max(0, n-M+1)}^{\min(n, L-1)} x[k] \cdot h[n - k], \quad \text{для } 0 \leq n \leq N-1$

\vspace{20pt}

Свойства дискретной свертки:
\begin{itemize}
    \item \textbf{Коммутативность:} \( x * h = h * x \)
    \item \textbf{Ассоциативность:} \( (x * h) * g = x * (h * g) \)
    \item \textbf{Дистрибутивность:} \( x * (h + g) = (x * h) + (x * g) \)
\end{itemize}

\vspace{20pt}

Передавать мы будем биты после свертки. Перед этим биты нужно удлинить, за это будет отвечать функция upsample, которая из одного бита делает напирмер 10.
То есть у нас было [1011] битов, а стало(если upsample до 5) [10000,00000,10000,10000], далее эти биты идут в функцию под названием filter, которая выполняет
свертку с еденичным импольсом после которой у нас получается такой результат [11111,00000,11111,11111]. Получается, что мы продлеваем наш сигнал.

\section*{BPSK и QPSK}

\textbf{BPSK} (Binary Phase Shift Keying) — это цифровая модуляция, при которой двоичные символы $0$ и $1$ передаются изменением фазы несущего сигнала на $0$ и $\pi$ радиан соответственно. Сигнал BPSK можно записать как:
\[
s_{\text{BPSK}}(t) = A \cos(2\pi f_c t + \phi), \quad \phi = 
\begin{cases}
0, & \text{для символа } 1, \\
\pi, & \text{для символа } 0.
\end{cases}
\]

\textbf{QPSK} (Quadrature Phase Shift Keying) — это модуляция, передающая два бита за символ с использованием четырёх фазовых состояний, разнесённых на $\frac{\pi}{2}$ радиан. Сигнал QPSK имеет вид:
\[
s_{\text{QPSK}}(t) = A \cos\!\left(2\pi f_c t + \frac{(2k + 1)\pi}{4}\right), \quad k = 0, 1, 2, 3,
\]
что соответствует фазам $\frac{\pi}{4}, \frac{3\pi}{4}, \frac{5\pi}{4}, \frac{7\pi}{4}$.

BPSK обеспечивает более высокую помехоустойчивость, тогда как QPSK удваивает спектральную эффективность по сравнению с BPSK при том же энергетическом расходе на бит.

\vspace{20pt}

\begin{figure}[H]
    \centering
    \includegraphics[width=0.8\textwidth]{qpskbpsk_vec.png}
    \caption{Векторная диаграмма BPSK, QPSK}
\end{figure}

\vspace{20pt}

На осцилографе QPSK сигнала, должно быть видно как меняется фаза

\begin{figure}[H]
    \centering
    \includegraphics[width=1.0\textwidth]{oscil_qpsk.png}
    \caption{Осциллограмма QPSK сигнала}
\end{figure}

\vspace{20pt}

Структурная схема QPSK модулятора будет выглядеть так:

\vspace{20pt}

\begin{figure}[H]
    \centering
    \includegraphics[width=0.8\textwidth]{qpsk_struct.png}
    \caption{Структурная схема QPSK модулятора}
\end{figure}

\section*{Выполнение}

Вместо BPSK, будем использовать QPSK. Напишем код для перевода в QPSK символы

\begin{lstlisting}[language=C++, caption=map для QPSK]
// QPSK символы
const complex<double> i0 (0, 0);
const complex<double> i1 (1, 1);
const complex<double> i2 (-1, 1);
const complex<double> i3 (-1, -1);
const complex<double> i4 (1, -1);

static map<string, complex<double>> qpsk_map = {
    {"00", i1},
    {"01", i2},
    {"11", i3},
    {"10", i4}
};
\end{lstlisting}

Функция для UpSampler, которая удлиняет сигнал, реализована с помощью простых циклов

\begin{lstlisting}[language=C++, caption=UpSampler]
void UpSampler(complex<double> *symbols, int len_s, complex<double> *symbols_ups, int L){
    for (size_t i = 0; i < (size_t)len_s*(size_t)L; i++) symbols_ups[i] = i0;
    for (size_t i = 0; i < (size_t)len_s; i++) symbols_ups[i*L] = symbols[i];
}
\end{lstlisting}

Функция для свертки, называется в моем коде filter, реализует процесс свертки массива уже после upsample с еденичным импульсом с заданной нами длиной

\begin{lstlisting}[language=C++, caption=filter]
void filter(complex<double> *symbols_ups, int len_symbols_ups, complex<double> *impulse, int L) {
vector<complex<double>> sum(len_symbols_ups, 0);
for (size_t i = 0; i < (size_t)len_symbols_ups; i++) {
    for (size_t j = 0; j < (size_t)L && (int)(i-j) >= 0; j++) {
        sum[i] += impulse[j] * symbols_ups[i-j];
    }
}
memcpy(symbols_ups, sum.data(), len_symbols_ups * sizeof(complex<double>));
}
\end{lstlisting}

Функция Mapper, которая формирует из 2 битов один символ, который далее соотносится к соотвующему комплексному значению в нашем qpsk\_map.

\begin{lstlisting}[language=C++, caption=Mapper]
void Mapper(int16_t *bits, int len_b, complex<double> *symbols){
    string pair_bits;
    for (size_t i = 0; i < (size_t)len_b; i += 2){
        pair_bits = to_string(bits[i]) + to_string(bits[i+1]);
        symbols[i/2] = qpsk_map[pair_bits];
    }
}
\end{lstlisting}

Для практики я рандомно формирую 20 битов от 0 до 1, затем пропускаю их через Mapper для того чтобы сформировать из битов символы по 2 бита каждый,
далее пропускаю эти символы через UpSampler чтобы "продлить" сигнал, затем подаю уже удлинненые символы на filter, который делает свертку с еденичным импульсом

\begin{lstlisting}[language=C++, caption=Как это устроено в коде]
int n = 20;
int16_t *bits = (int16_t*)malloc(n * sizeof(int16_t));
for (int i = 0; i < n; i++) bits[i] = rand() % 2;

int len_symbols = n / 2;
complex<double> *symbols = (complex<double>*)malloc(len_symbols * sizeof(complex<double>));
int L = 10;
int len_symbols_ups = len_symbols * L;
complex<double> *symbols_ups = (complex<double>*)malloc(len_symbols_ups * sizeof(complex<double>));
complex<double> impulse[L];
for (int i = 0; i < L; i++) impulse[i] = 1.0;

Mapper(bits, n, symbols);
UpSampler(symbols, len_symbols, symbols_ups, L);
filter(symbols_ups, len_symbols_ups, impulse, L);
// Show_Array("Символы", symbols_ups, len_symbols_ups);

int16_t *tx_samples = (int16_t*)malloc(2 * config.tx_mtu * sizeof(int16_t));
for (size_t i = 0; i < (size_t)config.tx_mtu; i++) {
    tx_samples[2*i] = (int16_t)((real(symbols_ups[i])) * 1000) << 4;
    tx_samples[2*i + 1] = (int16_t)((imag(symbols_ups[i])) * 1000) << 4;
}
    
// memcpy(config.tx_buff, tx_samples, 2 * len_symbols_ups * sizeof(int16_t));
fwrite(tx_samples, sizeof(int16_t), 2 * len_symbols_ups, tx);
\end{lstlisting}

Так же для того чтобы увидеть наши биты нам нужен код на python, который будет визулизировать все, что мы получили

\begin{lstlisting}[language=python, caption=Визуализация полученных битов]
import matplotlib.pyplot as plt
import numpy as np

print("""Choose 0 - Tx
Choose 1 - Rx""")

answer = int(input())

if answer == 1:
    nx = np.fromfile(f"/home/someotb/Files/Code/sdr/pluto/dev/build/rx.pcm", dtype=np.int16)
else:
    nx = np.fromfile(f"/home/someotb/Files/Code/sdr/pluto/dev/build/tx.pcm", dtype=np.int16)

rx = np.fromfile(f"/home/someotb/Files/Code/sdr/pluto/dev/build/rx.pcm", dtype=np.int16)
tx = np.fromfile(f"/home/someotb/Files/Code/sdr/pluto/dev/build/tx.pcm", dtype=np.int16)
print(f"RX: {rx}\nTX: {tx}")

samples = []

for x in range(0, len(nx), 2):
    samples.append(nx[x] + 1j * nx[x+1])

ampl = np.abs(samples)
phase = np.angle(samples)
time = np.arange(len(samples))

# plt.subplot(3,1,1)
# plt.plot(time, ampl)
# plt.grid(True)
# plt.tight_layout()
# plt.title("Amplitude")

# plt.subplot(3,1,2)
# plt.plot(time, phase)
# plt.grid(True)
# plt.tight_layout()
# plt.title("Phase")


# plt.subplot(3,1,3)
plt.plot(time, nx[0::2])
# plt.plot(time, nx[1::2])
plt.grid(True)
plt.tight_layout()
plt.title("Real and Imag parts")
plt.show()
\end{lstlisting}

\vspace{20pt}

По итогу мы можем считатать наши pcm файлы, в которых содержатся наши сгенерированные биты от 0 до 1, можем вывести их амплитуду, фазу и взять
реальную и мнимую части

\vspace{20pt}

\begin{figure}[H]
    \centering
    \includegraphics[width=0.8\textwidth]{result.png}
    \caption{Реальная и мнимая часим полученных битов}
\end{figure}

\section*{Вывод}

Получилось отправить и принять, сгенерированные биты, которые прошли через Mapper, UpSampler, filter, отобразить их амплитуду, фазу,
реальную и мнимую части, получить графики.